\subsection{H. WEYL 积分公式}

令 e 为 G 的单位元，\(\bar{e}\) 为其在 G/T 中的类。将 G 在 e 处的切空间与 g 识别，将 T 在 e 处的切空间与 t 识别，并将 G/T 在 \(\bar{e}\) 处的切空间与 g/t 识别。将 \((u,v)\mapsto u\cap v\) 记为下列 R-双线性映射

\[
\operatorname{Alt} ^ {n - r} (\mathfrak {g} / \mathfrak {t}) \times \operatorname{Alt} ^ {r} (\mathfrak {t}) \rightarrow \operatorname{Alt} ^ {n} (\mathfrak {g})
\]

它在第 1 节中已有定义。

回忆（第 III 章，第 3 节，第 13 号，命题 50），映射 \(\omega \mapsto \omega(e)\) 是从 G（分别为 T）上的 \(n\)（分别为 \(r\)）次左不变微分形式的向量空间到 \(\mathrm{Alt}^n(\mathfrak{g})\)（分别为 \(\mathrm{Alt}^r(t)\)）的同构。进一步注意，由于 \(\mathbf{R}^*\) 的每个连通紧子群都退化为单位元，\(\det \mathrm{Ad}g = 1\) 对所有 \(g \in \mathrm{G}\) 成立，所以 G 上的 \(n\) 次左不变微分形式同时也是右不变的，并且在内自同构下不变（第 III 章，第 3 节，第 16 号，命题 54 的推论）：从现在起我们简称为 G-不变微分形式。

同样，由第 III 章第 3 节第 16 号命题 56 及上述结果可知，映射 \(\omega \mapsto \omega (\bar{e})\) 是从次数为 \(n - r\) 的 \(\mathrm{G / T}\) 上的 G-不变微分形式空间到 \(\mathrm{Alt}^{n - r}(\mathfrak{g} / \mathfrak{t})\) 的同构。

若 \(\omega_{G/T}\) 是 G/T 上次数为 n - r 的 G-不变微分形式，而 \(\omega_{T}\) 是 T 上次数为 r 的不变微分形式，则将 G 上满足下式的唯一 n 次不变微分形式记为 \(\omega_{G/T} \cap \omega_{T}\)：

\[
(\omega_ {\mathrm{G/T}} \cap \omega_ {\mathrm{T}}) (e) = \omega_ {\mathrm{G/T}} (\bar {e}) \cap \omega_ {\mathrm{T}} (e).
\]

最后回忆，\(f : (\mathrm{G}/\mathrm{T}) \times \mathrm{T} \to \mathrm{G}\) 表示如下流形态射：它由映射 \((g, t) \mapsto gtg^{-1}\) 从 \(G \times T\) 到 G 通过取商诱导而来（第 5 节，第 4 号）。若 \(\alpha\) 和 \(\beta\) 分别是 G/T 和 T 上的微分形式，则将 \(\alpha \wedge \beta\) 简记为形式 \(pr_{1}^{*}\alpha \wedge pr_{2}^{*}\beta\) 在 \((\mathrm{G}/\mathrm{T}) \times \mathrm{T}\) 上的形式。

对于 \(t \in \mathrm{T}\)，将 \(\operatorname{Ad}_{\mathfrak{g} / \mathfrak{t}}(t)\) 记为 \(\mathfrak{g} / \mathfrak{t}\) 的自同态，它由 \(\operatorname{Ad} t\) 通过取商诱导而来。置

\[
\delta_ {\mathrm{G}} (t) = \det (\mathrm{Ad} _ {\mathfrak {g} / \mathfrak {t}} (t) - 1) = \prod_ {\alpha \in \mathrm{R} (\mathrm{G}, \mathrm{T})} (t ^ {\alpha} - 1).\tag{3}
\]

令 \(x \in \mathfrak{t}\) 且 \(\alpha \in \mathrm{R}(\mathrm{G},\mathrm{T})\)；将 \(\hat{\alpha}\) 记作元素 \((2\pi i)^{-1}\delta (\alpha)\)（它属于 \(\mathfrak{t}^*\)），于是

\[
((\exp x) ^ {\alpha} - 1) ((\exp x) ^ {- \alpha} - 1) = (e ^ {2 \pi i \hat {\alpha} (x)} - 1) (e ^ {- 2 \pi i \hat {\alpha} (x)} - 1) = 4 \sin^ {2} \pi \hat {\alpha} (x).
\]

若 \(\mathrm{R}_{+}(\mathrm{G},\mathrm{T})\) 表示相对于基 B 的 \(\mathrm{R}(\mathrm{G},\mathrm{T})\) 的正根集合，则有

\[
\delta_ {\mathrm{G}} (\exp x) = \prod_ {\alpha \in \mathrm{R} _ {+} (\mathrm{G}, \mathrm{T})} 4 \sin^ {2} \pi \hat {\alpha} (x),
\]

所以特别地，\(\delta_{\mathrm{G}}(t)>0\) 对所有 \(t\in T_{r}\) 成立。我们还注意到，\(\delta_{\mathrm{G}}(t)=\delta_{\mathrm{G}}(t^{-1})\) 对 \(t\in T\) 成立。

命题 1. 设 \(\omega_{\mathrm{G}}, \omega_{\mathrm{G/T}}\) 和 \(\omega_{\mathrm{T}}\) 分别是 G、G/T 和 T 上次数分别为 \(n, n - r\) 和 \(r\) 的不变微分形式。若 \(\omega_{\mathrm{G}} = \omega_{\mathrm{G/T}} \cap \omega_{\mathrm{T}}\)，则

\[
f ^ {*} (\omega_ {\mathrm{G}}) = \omega_ {\mathrm{G/T}} \wedge \delta_ {\mathrm{G}} \omega_ {\mathrm{T}}.
\]

显然可以假设 \(\omega_{G/T}\) 和 \(\omega_{T}\) 均非零；于是微分形式 \((u,t)\mapsto\omega_{\mathrm{G/T}}(u)\wedge\omega_{\mathrm{T}}(t)\) 在 \((\mathrm{G/T})\times\mathrm{T}\) 上是 n 次且处处非零的；因此存在 \(\delta\) 在 \((\mathrm{G/T})\times\mathrm{T}\) 上的数值函数，使得

\[
f ^ {*} (\omega_ {\mathrm{G}}) (u, t) = \delta (u, t) \omega_ {\mathrm{G/T}} (u) \wedge \omega_ {\mathrm{T}} (t).
\]

现在注意，对于 \(h \in \mathrm{G}\)、\(u \in \mathrm{G} / \mathrm{T}\)、\(t \in \mathrm{T}\)，有 \(f(h.u,t) = (\operatorname{Int} h)f(u,t)\)；由于 \(\omega_{\mathrm{G}}\) 在内自同构下不变，立即得到 \(\delta(h.u,t) = \delta(u,t)\)，从而 \(\delta(u,t) = \delta(\bar{e},t)\)。

将商映射 \(p:\mathfrak{g}\to \mathfrak{g} / \mathfrak{t}\) 记为 p，并将由下式定义的映射 \(\varphi :\mathfrak{g} / \mathfrak{t}\rightarrow \mathfrak{g}\) 记为：

\[
\varphi (p (X)) = (\operatorname{Ad} t ^ {- 1}) X - X \text { for } X \in \mathfrak {g};
\]

回忆（第 5 节，第 4 号，引理 4），切映射

\[
\mathrm{T} _ {(e, t)} (f): \mathrm{T} _ {e} (\mathrm{G} / \mathrm{T}) \times \mathrm{T} _ {t} (\mathrm{T}) \to \mathrm{T} _ {t} (\mathrm{G})
\]

将 \((z,tH)\) 送到 \(t(\varphi (z) + H)\)，其中 \(z\in \mathfrak{g} / \mathfrak{t},H\in \mathfrak{t}\)。

令 \(z_{1},\ldots,z_{n-r}\) 为 g/t 中的元素，令 \(H_{1},\ldots,H_{r}\) 为 t 中的元素。则

\[
\begin{array}{l} f ^ {*} \omega_ {\mathrm{G}} (\bar {e}, t) (z _ {1}, \ldots , z _ {n - r}, t H _ {1}, \ldots , t H _ {r}) \\ \quad = \omega_ {\mathrm{G}} (t) (t \varphi (z _ {1}), \ldots , t \varphi (z _ {n - r}), t H _ {1}, \ldots , t H _ {r}) \text {by the calculation of T} _ {(\bar {e}, t)} (f) \\ \quad = \omega_ {\mathrm{G}} (e) (\varphi (z _ {1}), \ldots , \varphi (z _ {n - r}), H _ {1}, \ldots , H _ {r}) \text {since} \omega_ {\mathrm{G}} \text {is invariant} \\ \quad = \omega_ {\mathrm{G/T}} (\bar {e}) (p \varphi (z _ {1}), \ldots , p \varphi (z _ {n - r})) \cdot \omega_ {\mathrm{T}} (e) (H _ {1}, \ldots , H _ {r}) \text {(no. 1, Lemma 1)} \\ \quad = \det (p \varphi) \omega_ {\mathrm{G/T}} (\bar {e}) (z _ {1}, \ldots , z _ {n - r}). \omega_ {\mathrm{T}} (e) (H _ {1}, \ldots , H _ {r}) \\ \quad = \delta_ {\mathrm{G}} (t) \omega_ {\mathrm{G/T}} (\bar {e}) (z _ {1}, \ldots , z _ {n - r}). \omega_ {\mathrm{T}} (t) (t H _ {1}, \ldots , t H _ {r}) \\ \qquad \qquad \qquad \qquad \qquad \qquad \text {since} \omega_ {\mathrm{T}} \text {is invariant} \\ \quad = \delta_ {\mathrm{G}} (t) (\omega_ {\mathrm{G/T}} \wedge \omega_ {\mathrm{T}}) (\bar {e}, t) (z _ {1}, \ldots , z _ {n - r}, t H _ {1}, \ldots , t H _ {r}), \end{array}
\]

所以 \(f^{*}\omega_{\mathrm{G}}(\bar{e},t) = \delta_{\mathrm{G}}(t)(\omega_{\mathrm{G / T}}\wedge \omega_{\mathrm{T}})(\bar{e},t)\)；于是 \(\delta (\bar {e},t) = \delta_{\mathrm{G}}(t)\)，命题得证。

分别以形式 \(\omega_{G}, \omega_{T}\) 和 \(\omega_{G/T}\) 所定义的定向赋予流形 G、T 和 G/T。这些形式在 G、T 和 G/T 上定义不变测度（第 III 章，第 3 节，第 16 号，命题 55 和 56），这些测度也记为 \(\omega_{G}, \omega_{T}\) 和 \(\omega_{G/T}\)。

引理 2. 若 \(\omega_{\mathrm{G}} = \omega_{\mathrm{G / T}}\cap \omega_{\mathrm{T}}\)，则

\[
\int_ {\mathrm{G}} \omega_ {\mathrm{G}} = \int_ {\mathrm{G/T}} \omega_ {\mathrm{G/T}}. \int_ {\mathrm{T}} \omega_ {\mathrm{T}}.
\]

将从 G 到 G/T 的典则态射记为 \(\pi\)。令 \(g \in G\)，并令 \(t_{1}, \ldots, t_{n-r}\) 为 \(\mathrm{T}_{\pi(g)}(\mathrm{G}/\mathrm{T})\) 中的元素。将纤维 \(\pi^{-1}(\pi(g)) = g\mathrm{T}\) 通过平移 \(\gamma(g)\) 与 T 识别。关系 \(\omega_{G} = \omega_{G/T} \cap \omega_{T}\) 于是推出下列等式（《微分与解析流形，结果》，11.4.5）：

\[
\omega_ {\mathrm{G} \sqcup} (t _ {1}, \dots , t _ {n - r}) = (\omega_ {\mathrm{G} / \mathrm{T}} (t _ {1}, \dots , t _ {n - r})) \omega_ {\mathrm{T}}.
\]

因此 \(\int_{\pi}\omega_{\mathrm{G}} = \left(\int_{\mathrm{T}}\omega_{\mathrm{T}}\right)\omega_{\mathrm{G / T}}\)（《微分与解析流形，结果》，11.4.6），并且

\[
\int_ {\mathrm{G}} \omega_ {\mathrm{G}} = \int_ {\mathrm{G/T}} \int_ {\pi} \omega_ {\mathrm{G}} = \int_ {\mathrm{T}} \omega_ {\mathrm{T}}. \int_ {\mathrm{G/T}} \omega_ {\mathrm{G/T}}
\]

（《微分与解析流形，结果》，11.4.8）。

引理 3. 在 \((\mathrm{G}/\mathrm{T}) \times \mathrm{T}_{r}\) 上，\(G_{r}\) 上测度 dg 在局部同胚 \(f_{r}\) 下的逆像是测度 \(\mu \otimes \delta_{G} dt\)（积分，第 V 章，第 6 节，第 6 号），其中 \(\mu\) 是 G/T 上总质量为 1 的唯一 G-不变测度。

在 T（分别为 G/T）上选取最高次数的不变微分形式 \(\omega_{T}\)（分别为 \(\omega_{G/T}\)），使得 \(\omega_{T}\)（分别为 \(\omega_{G/T}\)）所定义的测度等于 dt（分别为 \(\mu\)）。置 \(\omega_{G} = \omega_{G/T} \cap \omega_{T}\)。引理 2 蕴含，\(\omega_{G}\) 所定义的测度等于 dg。令 U 为 (G/T) × T \(_{r}\) 的开子集，使得 \(f_{r}\) 将 U 同构地映到 G \(_{r}\) 的开子集 V。令 \(\varphi\) 为 V 中具有紧支集的连续函数；仍以 \(\varphi\) 表示 \(\varphi\) 延拓到 G \(_{r}\) 且在 V 外为零的函数。有

\[
\begin{array}{r l} & {\int_ {\mathrm{V}} \varphi d g = \int_ {\mathrm{V}} \varphi \omega_ {\mathrm{G}} = \int_ {\mathrm{U}} (\varphi \circ f _ {r}) f _ {r} ^ {*} (\omega_ {\mathrm{G}})} \\ & {\qquad = \int_ {\mathrm{U}} (\varphi \circ f _ {r}) \omega_ {\mathrm{G/T}} \wedge \delta_ {\mathrm{G}} \omega_ {\mathrm{T}} \quad (\text {Prop. 1})} \\ & {\qquad = \int_ {\mathrm{U}} (\varphi \circ f _ {r}) d \mu . \delta_ {\mathrm{G}} d t,} \end{array}
\]

引理得证。

定理 1（H. Weyl）。测度 \(dg\)（定义在 \(G\) 上）是映射 \((g, t) \mapsto gtg^{-1}\) 从 \(G \times T\) 到 \(G\) 下测度 \(dg \otimes \frac{1}{w(G)}\delta_G dt\) 的像，其中

\[
\delta_ {\mathrm{G}} (t) = \det (\operatorname{Ad} _ {\mathfrak {g} / \mathfrak {t}} (t) - 1) = \prod_ {\alpha \in \mathrm{R} (\mathrm{G}, \mathrm{T})} (t ^ {\alpha} - 1).
\]

等价地（积分，第 V 章，第 6 节，第 3 号，命题 4），\(dg\) 是映射 \(f: (\mathrm{G} / \mathrm{T}) \times \mathrm{T} \to \mathrm{G}\) 下测度 \(\mu \otimes \frac{1}{w(\mathrm{G})}\delta_{\mathrm{G}}dt\) 的像。

我们证明最后一个断言。由第 5 节第 1 号以及《微分与解析流形，结果》10.1.3 c) 可知，G-G \(_{r}\) 在 G 中是零测集，T-T \(_{r}\) 在 T 中也是零测集。此外，映射 \(f_{r}\) 使 (G/T) × T \(_{r}\) 成为 G \(_{r}\) 的主覆盖，其群为 W（第 5 节，第 4 号，命题 4 b)）。定理现由引理 3 以及积分，第 V 章第 6 节第 6 号命题 11 得出。

推论 1. (i) 令 \(\varphi\) 为 G 上取值于 Banach 空间或 \(\overline{\mathbf{R}}\) 的可积函数。对于几乎所有 \(t \in \mathrm{T}\)，函数 \(g \mapsto \varphi(gtg^{-1})\) 在 G 上关于 \(dg\) 可积。函数 \(t \mapsto \delta_{\mathrm{G}}(t) \int_{\mathrm{G}} \varphi(gtg^{-1}) dg\) 在 T 上可积，并且

\[
\int_ {\mathrm{G}} \varphi (g) d g = \frac {1}{w (\mathrm{G})} \int_ {\mathrm{T}} \left(\int_ {\mathrm{G}} \varphi (g t g ^ {- 1}) d g\right) \delta_ {\mathrm{G}} (t) d t\tag{4}
\]

（“Hermann Weyl 积分公式”。）

\begin{enumerate}
\def\labelenumi{(\roman{enumi})}
\setcounter{enumi}{1}
\tightlist
\item
  令 \(\varphi\) 为 \(\mathbf{G}\) 上的正可测函数。对于几乎所有 \(t \in \mathrm{T}\)，函数 \(g \mapsto \varphi(gtg^{-1})\) 在 \(\mathbf{G}\) 上可测。函数 \(t \mapsto \int_{\mathbf{G}}^{*} \varphi(gtg^{-1}) dg\) 在 \(\mathrm{T}\) 上可测，并且
\end{enumerate}

\[
\int_ {\mathrm{G}} ^ {*} \varphi (g) d g = \frac {1}{w (\mathrm{G})} \int_ {\mathrm{T}} ^ {*} \left(\int_ {\mathrm{G}} ^ {*} \varphi (g t g ^ {- 1}) d g\right) \delta_ {\mathrm{G}} (t) d t.\tag{5}
\]

由于态射 f 是由映射 \((g, t) \mapsto gtg^{-1}\) 从 \(G \times T \to G\) 通过取商诱导的，只需应用《积分》第 V 章第 5 节、第 6 节、第 8 节以及《积分》第 VII 章第 2 节。

推论 2. 令 \(\varphi\) 为 \(\mathbf{G}\) 上的中心函数（即满足 \(\varphi(gh) = \varphi(hg)\) 对所有 \(g\) 和 \(h\) 属于 \(\mathbf{G}\) 成立），其值属于 Banach 空间或 \(\overline{\mathbf{R}}\)。

\begin{enumerate}
\def\labelenumi{\alph{enumi})}
\item
  \(\varphi\) 可测，当且仅当其在 T 上的限制可测。
\item
  \(\varphi\) 可积，当且仅当函数 \((\varphi |\mathrm{T})\delta_{\mathrm{G}}\) 在 \(\mathrm{T}\) 上可积；在这种情况下，我们有
\end{enumerate}

\[
\int_ {\mathrm{G}} \varphi (g) d g = \frac {1}{w (\mathrm{G})} \int_ {\mathrm{T}} \varphi (t) \delta_ {\mathrm{G}} (t) d t.\tag{6}
\]

将 \(p: \mathrm{G} / \mathrm{T} \times \mathrm{T} \to \mathrm{T}\) 记为第二投影。我们有 \(\varphi \circ f = (\varphi | \mathrm{T}) \circ p\)；此外，\(p\) 下测度 \(\mu \otimes \frac{1}{w(\mathrm{G})} \delta_{\mathrm{G}} dt\) 的像为 \(\frac{1}{w(\mathrm{G})} \delta_{\mathrm{G}} dt\)。推论现由上面的定理 1 以及《积分》第 V 章第 6 节第 2 号定理 1 应用于两个真映射 \(f\) 和 \(p\) 后得到。

推论 3. 令 H 为含 T 的 G 的连通闭子群，h 为其李代数，dh 为 H 上总质量为 1 的 Haar 测度。令 \(\varphi\) 为 G 上取值于 Banach 空间或 R 的可积中心函数。于是函数 \(h \mapsto \varphi(h) \det(\mathrm{Ad}_{\mathfrak{g}/\mathfrak{h}}(h) - 1)\) 在 H 上可积且为中心函数，并且

\[
\int_ {\mathrm{G}} \varphi (g) d g = \frac {w (\mathrm{H})}{w (\mathrm{G})} \int_ {\mathrm{H}} \varphi (h) \mathrm{det} (\mathrm{Ad} _ {\mathfrak {g} / \mathfrak {h}} (h) - 1) d h.\tag{7}
\]

事实上，函数 \(h \mapsto \varphi(h) \det(\mathrm{Ad}_{\mathfrak{g}/\mathfrak{h}}(h) - 1)\) 是 H 上的中心函数，其在 T 上的限制为函数 \(t \mapsto \varphi(t) \delta_{\mathrm{G}}(t) \delta_{\mathrm{H}}(t)^{-1}\)。因此，推论由将推论 2 应用于 G 和 H 得到。

备注。1) 若在推论 3 中取 \(\varphi = 1\)，则得到

\[
\int_ {\mathrm{H}} \det (\mathrm{Ad} _ {\mathfrak {g} / \mathfrak {h}} (h) - 1) d h = w (\mathrm{G}) / w (\mathrm{H})\tag{8}
\]

并且特别有

\[
\int_ {\mathrm{T}} \delta_ {\mathrm{G}} (t) d t = w (\mathrm{G}).\tag{9}
\]

\begin{enumerate}
\def\labelenumi{\arabic{enumi})}
\setcounter{enumi}{1}
\tightlist
\item
  令 \(\nu\) 为商空间 T/W 上由下式定义的测度
\end{enumerate}

\[
\int_ {\mathrm{T/W}} \psi (\tau) d \nu (\tau) = \frac {1}{w (\mathrm{G})} \int_ {\mathrm{T}} \psi (\pi (t)) \delta_ {\mathrm{G}} (t) d t,
\]

其中 \(\pi\) 表示 T 到 T/W 的典则投影。推论 2 意味着，同胚 \(T / W \to G / \operatorname{Int}(G)\)（第 2 节，第 5 号，命题 5 的推论 1）

将测度 \(\nu\) 传送为典则投影 \(\mathrm{G} \rightarrow \mathrm{G}/\mathrm{Int}(\mathrm{G})\) 下测度 dg 的像。

\begin{enumerate}
\def\labelenumi{\arabic{enumi})}
\setcounter{enumi}{2}
\tightlist
\item
  假设 G 单连通。令 A 为 t 的一个单房，令 dx 为 t 上满足 \(\int_{A} dx = 1\) 的 Haar 测度。则测度 \(\nu\) 也可由测度 \(\frac{1}{w(G)} \prod_{\alpha \in R_{+}(G,T)} 4 \sin^{2} \pi \hat{\alpha}(x) dx\) 在 \(\overline{A}\) 上通过同胚 \(\overline{A} \rightarrow T/W\) 传送而得到（第 5 节，第 2 号，命题 2 的推论 1）。
\end{enumerate}

例. 取 G 为群 \(\mathbf{SU}(2,\mathbf{C})\)，取 T 为对角矩阵子群（\(\S3\)，第 6 号）；通过选取 t 的基 \(\{iH\}\)（同上引文）将 t 与 R 识别。置 A = 10, \(\pi\)；这是 t 的一个单房。区间 \(\overline{A} = [0, \pi]\) 可与 G 的共轭类空间识别，其中元素 \(\theta\) 属于 \(\overline{A}\)，对应于矩阵 \(\begin{pmatrix} e^{i\theta} & 0 \\ 0 & e^{-i\theta} \end{pmatrix}\) 所属的共轭类。令 \(d\theta\) 为 \([0, \pi]\) 上的 Lebesgue 测度；由上述结果可知，G 上的 Haar 测度在 \(\overline{A}\) 上的像是测度 \(\frac{2}{\pi}\sin^{2}\theta d\theta\)。

